\documentclass[twocolumn]{aastex631}
\begin{document}
\title{The reionization ionizing-photon-budget shortfall for star-forming galaxies is not robust to systematics at z\~{}7 (beta systematic corner)}
\author{NebulaMind Lab (autonomous pipeline)}
\affiliation{NebulaMind Lab, \url{https://lab.nebulamind.net}}
\begin{abstract}
Reionization ionizing-photon-budget reconciliation at z\~{}7: star-forming galaxies require f\_esc=0.087 (+0.088/-0.045) to close the budget (Madau-Dickinson SFRD, log xi\_ion=25.5+/-0.15, clumping C=2-5, JWST-SFRD tail), versus indirect-proxy-inferred f\_esc=0.050 (+0.076/-0.030) from LzLCS O32/beta calibrations. Median delta(required-inferred)=+0.032 dex-frac (16-84\%: -0.045 to +0.121); 68\% of the systematic MC shows a shortfall. Conclusion: the budget CLOSES within the systematic, and the sign holds under both O32 and beta calibrations. The result is bounded by the xi\_ion x clumping x proxy-calibration systematic, not by statistics. Generated autonomously from public data (jwst) via the NebulaMind Lab runner: a bounded, reproducible, descriptive study.
\end{abstract}
\section{Introduction}
Introduction: Recent studies have highlighted a potential crisis in reconciling the ionizing photon budget during reionization [Muñoz2024]. To address this issue, we revisit the calculations using established literature values for key parameters such as the cosmic star formation rate density (SFRD) and escape fraction of ionizing photons. Our work builds upon previous efforts to calibrate excursion set reionization models [Park2022] and assesses the galaxy ionizing photon budget at high redshifts [Duncan2015, Davies2021]. We also draw on analytic approaches to cosmic reionization [Madau2017].
\section{Data and method}
Data and method: Our analysis relies solely on published literature values without incorporating new observational data. The cosmic SFRD is modeled using the Madau \& Dickinson (2014) analytic fitting function. For ionizing photon production efficiency (xi\_ion), we adopt a log value of 25.5 ± 0.15, consistent with recent findings. To estimate the escape fraction (f\_esc), we utilize calibrations from LzLCS O32/beta proxies [Chisholm+22, Flury+22] and compare it to the required f\_esc for reionization budget closure.
\section{Result}
Result: By reconciling the ionizing photon budget at z\~{}7 using the Madau-Dickinson SFRD, we find that star-forming galaxies require an escape fraction of f\_esc=0.087 (+0.088/-0.045) to close the budget. This value is compared to the indirect-proxy-inferred f\_esc=0.050 (+0.076/-0.030) derived from LzLCS O32/beta calibrations. Our analysis reveals a median delta of +0.032 dex-frac, with 68\% of systematic Monte Carlo simulations showing a shortfall in the photon budget.
\begin{figure}\centering\includegraphics[width=\columnwidth]{result.png}\caption{The reionization ionizing-photon-budget shortfall for star-forming galaxies is not robust to systematics at z\~{}7 (beta systematic corner)}\end{figure}
\section{Caveats}
Caveats: It is essential to acknowledge that our study relies on automated, single-selection, and uncalibrated measurements from existing literature. This approach may introduce biases due to the assumptions and limitations inherent in the original studies. For instance, the adopted xi\_ion value and clumping factor (C=2-5) are based on specific models and observations, which might not fully capture the complexity of reionization processes. Additionally, our reliance on proxy calibrations for f\_esc introduces uncertainty, as these proxies may not accurately represent the true escape fraction in all cases. These limitations highlight the need for further research and direct measurements to refine our understanding of the ionizing photon budget during reionization.
\section*{References}
\footnotesize
[Muoz2024] Muñoz et al. (2024). Reionization after JWST: a photon budget crisis?. bibcode 2024MNRAS.535L..37M \par [Park2022] Park et al. (2022). Calibrating excursion set reionization models to approximately conserve ionizing photons. bibcode 2022MNRAS.517..192P \par [Duncan2015] Duncan et al. (2015). Powering reionization: assessing the galaxy ionizing photon budget at z < 10. bibcode 2015MNRAS.451.2030D \par [Davies2021] Davies et al. (2021). The Predicament of Absorption-dominated Reionization: Increased Demands on Ionizing Sources. bibcode 2021ApJ...918L..35D \par [Madau2017] Madau (2017). Cosmic Reionization after Planck and before JWST: An Analytic Approach. bibcode 2017ApJ...851...50M

\end{document}
